% Vietnamese translation for OI Public Library
% Source-PDF: 图论 GRAPH THEORY/Making Graphs into Trees - immortalCO, WrongAnswer.pdf
\documentclass[11pt]{article}
\usepackage{oipl-vn}
\usepackage{tikz}

\newcommand{\VT}{V_T}
\newcommand{\ET}{E_T}
\newcommand{\RT}{R_T}
\newcommand{\ST}{S_T}

\title{Biến đồ thị thành cây}
\author{make\_pair(immortalCO, WrongAnswer)}
\date{11 tháng 3 năm 2017}

\begin{document}
\maketitle

\origtitle{Making Graphs into Trees}
\sourcepath{Graph Theory / Making Graphs into Trees - immortalCO, WrongAnswer}

\begin{abstract}
Trong nhiều bài toán đồ thị của OI, nếu có thể biến đồ thị thành một cấu trúc
dạng cây thì lời giải trở nên thuận tiện hơn nhiều. Tài liệu gốc tổng kết ba
cách nhìn chính: cây DFS, cây BFS và cây tròn-vuông của đồ thị cactus. Bản
dịch chuyển bộ slide thành ghi chú bài giảng LaTeX, giữ lại các định nghĩa,
tính chất và ứng dụng thuật toán chính; các hình minh họa được diễn giải bằng
lời khi chúng chỉ phục vụ ví dụ hoặc thứ tự duyệt.
\end{abstract}

\tableofcontents

\section{Cây DFS}

Cho đồ thị liên thông vô hướng \(G=(V,E)\). Khi chạy DFS từ một gốc \(r\), các
cạnh được DFS đi qua tạo thành một cây
\[
  T=(V,E').
\]
Ta gọi các cạnh trong \(E'\) là \term{cạnh cây}, còn các cạnh trong
\(E\setminus E'\) là \term{cạnh ngoài cây}.

Trong cây DFS của đồ thị vô hướng, mọi cạnh ngoài cây \((u,v)\) đều nối một
đỉnh với một tổ tiên của nó trong \(T\). Với cạnh ngoài cây \((u,v)\), ta nói
nó \term{phủ} các cạnh cây nằm trên đường đi giữa \(u\) và \(v\) trong \(T\).
Tính chất đơn giản này là lý do cây DFS thường biến các ràng buộc chu trình
thành ràng buộc trên đường đi trong cây.

\subsection{Xử lý cactus bằng cây DFS}

Một cactus là đồ thị vô hướng liên thông trong đó mỗi cạnh thuộc nhiều nhất
một chu trình đơn. Sau khi dựng cây DFS của cactus, mỗi cạnh ngoài cây tương
ứng đúng với một chu trình: chu trình gồm chính cạnh ngoài cây đó và các cạnh
cây mà nó phủ. Vì các chu trình của cactus không giao nhau theo cạnh, hai cạnh
ngoài cây khác nhau sẽ phủ hai tập cạnh cây rời nhau. Nhờ đó ta có thể tìm
thuận tiện mỗi cạnh thuộc chu trình nào.

\subsection{Tập độc lập lớn nhất trên cactus}

Bài toán: cho một cactus, hãy tìm kích thước tập độc lập lớn nhất.

Cách trực tiếp bắt đầu từ cây DFS. Nếu bỏ qua các cạnh ngoài cây, ta có bài
toán tập độc lập trên cây với trạng thái
\[
  f(i,j),
\]
trong đó \(j\in\{0,1\}\) cho biết cha của \(i\) có được chọn hay không. Tuy
nhiên tập độc lập của cây DFS chưa chắc là tập độc lập của cactus, vì còn phải
thỏa các cạnh ngoài cây.

Do mỗi cạnh ngoài cây nằm trên một chu trình riêng, ta chỉ cần nhớ thêm trạng
thái của đỉnh trên cùng của chu trình chứa \(i\). Gọi \(k\in\{0,1\}\) là trạng
thái của đỉnh có độ sâu nhỏ nhất trên chu trình đó, và đặt
\[
  f(i,j,k)
\]
là kích thước tập độc lập lớn nhất trong cây con gốc \(i\), khi cha của \(i\)
có trạng thái \(j\) và đỉnh trên cùng của chu trình có trạng thái \(k\). Khi
\(i\) là đỉnh đáy của chu trình và \(k=1\), ta cấm chọn \(i\). Nhờ vậy các
ràng buộc ngoài cây được đưa về trạng thái cục bộ, và độ phức tạp vẫn tuyến
tính theo \(|V|\) hoặc \(|E|\), vốn cùng bậc trong cactus.

\subsection{Quy hoạch động trạng thái nén trên cây DFS}

Cây DFS cũng hữu ích khi đồ thị gần là cây, tức \(|E|-|V|\) nhỏ.

\paragraph{Đếm tô màu.}
Cho đồ thị liên thông vô hướng \(G=(V,E)\), cần đếm số cách tô \(V\) bằng
\(k\) màu sao cho hai đầu mỗi cạnh có màu khác nhau, lấy modulo \(p\). Giả sử
\[
  |V|\le 100,\qquad |E|<|V|+10.
\]
Nếu \(G\) là cây với \(n=|V|\), đáp án là
\[
  k(k-1)^{n-1}.
\]
Với đồ thị có ít cạnh ngoài cây, dựng cây DFS. Vì số cạnh ngoài cây không quá
\(10\), có thể nén trạng thái bằng cách ghi nhận quan hệ màu giữa các đầu mút
phía trên của các cạnh ngoài cây. Bản chất là xét các phân hoạch của tối đa
\(10\) phần tử; số phân hoạch là
\[
  \operatorname{Bell}(10)=115975.
\]
Một cách cài đặt dễ hơn là liệt kê bên ngoài phân hoạch màu của các đầu mút
đặc biệt, rồi bên trong chạy quy hoạch động trên cây.

Ta còn có thể giảm kích thước đồ thị trước khi DP. Với \(n>1\), mỗi lá trong
cây đóng góp nhân tử \(k-1\), nên có thể xóa lá lặp lại. Sau đó nén các chuỗi
đỉnh bậc \(2\) thành một cạnh có trọng số. Với một đường dài \(l\), tiền xử lý
hai giá trị \(a_l,b_l\): số cách tô phần trong đường khi hai đầu cùng màu và
khi hai đầu khác màu. Sau khi nén, mỗi cạnh mang hai trọng số tương ứng, và
DP chỉ còn chạy trên phần lõi của đồ thị. Nếu mọi đỉnh còn lại có bậc ít nhất
\(3\), từ \(\sum_v \deg(v)=2m\ge 3n\) suy ra \(n\le 2(m-n)\), nên độ phức tạp
phụ thuộc chủ yếu vào số dư chu trình \(m-n\).

\paragraph{Đếm tập độc lập kích thước \(k\).}
Biến thể khác là đếm số tập độc lập kích thước \(k\) modulo \(p\), khi
\(|E|-|V|\) nhỏ. Sau khi dựng cây DFS \(T\), đặt
\[
  f(i,j,S)
\]
là số cách chọn tập độc lập kích thước \(j\) trong cây con gốc \(i\). Tập
\(S\) lưu các ràng buộc do những cạnh ngoài cây phủ cạnh nối \(i\) với cha
của nó: cụ thể là cha của \(i\) và các đầu mút phía trên liên quan có được
chọn hay không. Nếu \(C_i\) là tập cạnh ngoài cây phủ cạnh cha của \(i\), độ
phức tạp có dạng
\[
  O\!\left(n^2\sum_{v\in V}2^{|C_v|}\right).
\]
Cách biểu diễn này vẫn nhanh ngay cả khi \(|E|-|V|\) không quá nhỏ, miễn là
mỗi cạnh cây chỉ bị ít cạnh ngoài cây phủ.

\section{Cây BFS}

Bên cạnh cây DFS còn có cây BFS. Tính chất của nó không thuận lợi cho DP cây
kiểu trên, vì cạnh ngoài cây không nhất thiết nối tổ tiên--hậu duệ. Tuy nhiên
nó có tính chất phân lớp rất mạnh: nếu chia đỉnh theo độ sâu trong cây BFS,
thì mọi cạnh ngoài cây chỉ nối hai đỉnh cùng lớp hoặc hai lớp kề nhau. Tư
tưởng này có thể dùng trong một số quy hoạch động trên đồ thị phân lớp.

Slide gốc khuyến nghị bài \emph{UOJ \#259}: đây là một bài nộp đáp án do tác
giả tự ra, và các kỹ thuật trong buổi nói chuyện có thể giải được khá nhiều
test. Gợi ý được ghi rõ là không nhất thiết phải lấy đỉnh \(1\) làm gốc.

\section{Cây tròn-vuông của cactus}

\subsection{Định nghĩa}

Một cactus là đồ thị vô hướng liên thông mà mỗi cạnh thuộc không quá một chu
trình đơn.

Với cactus \(G=(V,E)\), \term{cây tròn-vuông} của nó là đồ thị
\[
  T=(\VT,\ET)
\]
thỏa các điều kiện sau. Tập đỉnh được chia thành
\[
  \VT=\RT\cup\ST,\qquad \RT=V,\qquad \RT\cap\ST=\varnothing.
\]
Các đỉnh trong \(\RT\) gọi là \term{đỉnh tròn}; chúng tương ứng với đỉnh gốc
của cactus. Các đỉnh trong \(\ST\) gọi là \term{đỉnh vuông}; mỗi đỉnh vuông
tương ứng với một chu trình đơn của cactus.

Nếu một cạnh \(e\in E\) không thuộc chu trình nào, ta giữ cạnh đó trong
\(\ET\) như một cạnh nối hai đỉnh tròn. Với mỗi chu trình đơn \(R\), tạo một
đỉnh vuông \(p_R\), rồi nối \(p_R\) với mọi đỉnh tròn nằm trên \(R\).

\paragraph{Vì sao đây là cây?}
Các cạnh không thuộc chu trình được giữ nguyên nên không phá liên thông. Mỗi
chu trình được thay bằng một đỉnh vuông nối tới mọi đỉnh trên chu trình, nên
tính liên thông vẫn được bảo toàn. Nếu cactus có \(c=|E|-|V|+1\) chu trình,
thì
\[
  |\VT|=|V|+c=|E|+1.
\]
Mặt khác, một chu trình dài \(r\) có \(r\) cạnh trong cactus và cũng đóng góp
\(r\) cạnh từ đỉnh vuông tới các đỉnh tròn trong cây tròn-vuông, nên tổng số
cạnh vẫn là \(|E|\). Vì đồ thị thu được liên thông và có
\(|\VT|=|\ET|+1\), nó là một cây.

\subsection{Xây dựng}

Có thể xây dựng cây tròn-vuông bằng thuật toán Tarjan tìm thành phần song liên
thông đỉnh:
\begin{enumerate}
  \item chạy Tarjan từ một đỉnh bất kỳ;
  \item với mỗi thành phần song liên thông là một chu trình, lấy các đỉnh khỏi
        ngăn xếp; thứ tự trong ngăn xếp chính là thứ tự trên chu trình, tạo
        đỉnh vuông và nối tới các đỉnh tròn tương ứng;
  \item nếu một cạnh là cạnh cây không bị chu trình nào phủ, thêm trực tiếp
        cạnh tròn--tròn đó vào cây tròn-vuông.
\end{enumerate}

Một số tính chất cơ bản:
\begin{itemize}
  \item không có cạnh nối hai đỉnh vuông;
  \item cây tròn-vuông không phụ thuộc vào gốc DFS, ngoại trừ cách đánh số
        các đỉnh vuông;
  \item nếu đặt gốc \(r\), thì \term{sub-cactus} của một đỉnh tròn \(p\) là
        các đỉnh tròn nằm trong cây con của \(p\) trên cây tròn-vuông.
\end{itemize}

Ví dụ hình trong slide là mẫu 1 của bài \emph{UOJ 189: Tài xế tàu ra Tần
Xuyên}. Hình này dùng các đỉnh đánh số \(1,\ldots,14\) để minh họa trực quan
rằng một cactus quen thuộc có thể được thay bằng cây tròn-vuông: mỗi chu trình
được gom thành một đỉnh vuông, còn các cầu giữ nguyên là cạnh tròn--tròn.

\section{Ứng dụng của cây tròn-vuông}

\subsection{DP trên cactus}

\paragraph{BZOJ 4316: tập độc lập lớn nhất.}
Ràng buộc của slide là \(n\le 1000000,\ m\le 2000000\).
Ta dùng trạng thái giống cây:
\[
  f(i,0/1)
\]
là tập độc lập lớn nhất trong sub-cactus của \(i\), tùy \(i\) có được chọn hay
không. Nếu cạnh là tròn--tròn, chuyển giống cây. Nếu gặp cạnh tròn--vuông,
lấy toàn bộ các đỉnh trên chu trình tương ứng và chạy DP tập độc lập trên
vòng. Cách này có ít chi tiết hơn so với xử lý trực tiếp bằng cây DFS, vì
Tarjan đưa các đỉnh trên chu trình ra theo đúng thứ tự.

\paragraph{BZOJ 1023: đường kính cactus.}
Cần tìm giá trị lớn nhất của khoảng cách ngắn nhất giữa hai đỉnh. Với đỉnh
tròn, xử lý như đường kính cây: lưu độ sâu lớn nhất và nhì đi xuống. Với đỉnh
vuông, cần xét đường đi ngắn nhất dài nhất có LCA là chính chu trình đó. Trên
một chu trình, hai hướng đi tạo hai ứng viên; có thể dùng hàng đợi đơn điệu để
xử lý các giá trị cực đại trên vòng. Ràng buộc trong slide cũng là
\(n\le 1000000,\ m\le 2000000\).

\subsection{Truy vấn đường đi ngắn nhất}

Bài BZOJ 2125 yêu cầu nhiều truy vấn khoảng cách ngắn nhất giữa hai đỉnh trong
cactus, với \(n,m,q\le 10^6\). Hai đỉnh trong cactus có các đường đi đơn tạo
thành một chuỗi gồm một số cạnh cây và một số chu trình; trên mỗi chu trình,
đường ngắn nhất luôn chọn phía ngắn hơn.

Ví dụ cây tròn--vuông có trọng số minh họa đúng quy tắc này: cạnh tròn--tròn
giữ trọng số gốc, cạnh nối đỉnh
vuông với cha của nó có trọng số \(0\), còn các cạnh tròn--vuông khác nhận độ
dài ngắn nhất từ đỉnh tròn đó về cha của chu trình. Nhờ vậy, mọi chu trình đã
được ``quyết định'' đi phía ngắn hơn, trừ chu trình chứa LCA của truy vấn.

Gán trọng số cho cây tròn-vuông như sau. Chọn một đỉnh tròn làm gốc. Mọi cạnh
tròn--tròn giữ trọng số như trong đồ thị gốc. Với một cạnh tròn--vuông, nếu đó
là cạnh nối đỉnh vuông với cha của nó thì đặt trọng số \(0\); nếu không, đặt
trọng số bằng khoảng cách ngắn nhất từ đỉnh tròn đó đến cha của đỉnh vuông
trên chu trình.

Khi đó nếu LCA của hai đỉnh là đỉnh tròn, khoảng cách ngắn nhất chính là khoảng
cách có trọng số trên cây tròn-vuông. Nếu LCA là đỉnh vuông, còn phải xét riêng
chu trình của đỉnh vuông: dùng HLD hoặc bảng nhảy để tìm hai nhánh của truy
vấn trên chu trình, rồi cộng phía ngắn hơn giữa hai vị trí đó.

\subsection{Virtual cactus}

Tương tự virtual tree, ta có thể dựng \term{virtual cactus}. Trước hết cần một
tính chất: nếu một cây \(T=(\ST+\RT,\ET)\) thỏa không có cạnh nối hai đỉnh
vuông, thì nó là một cây tròn-vuông hợp lệ của một cactus nào đó. Chứng minh
có tính xây dựng: với mỗi đỉnh vuông, nối các đỉnh tròn kề nó thành một chu
trình.

Nếu trực tiếp lấy virtual tree của cây tròn-vuông, có thể xuất hiện cạnh ảo
nối hai đỉnh vuông. Để tránh điều đó, với mỗi cạnh ảo đi ra từ một đỉnh vuông,
chèn thêm một đỉnh tròn biểu diễn vị trí trên chu trình mà nhánh đó rời khỏi.
Kích thước vẫn là \(O(|S|)\) với \(S\) là tập đỉnh quan trọng. Cây ảo thu được
lại là một cây tròn-vuông hợp lệ, tương ứng với một virtual cactus.

Với tập điểm quan trọng \(\{2,8\}\), nếu lấy virtual tree trực tiếp trên cây
tròn-vuông thì cạnh ảo có thể nối hai đỉnh vuông. Sau khi chèn thêm đỉnh tròn
biểu diễn vị trí rời khỏi chu trình, ta lại thu được một virtual cây
tròn-vuông hợp lệ. Sơ đồ nhỏ dưới đây minh họa đúng điểm phải sửa.

\begin{center}
\begin{tikzpicture}[
  round/.style={circle,draw,minimum size=0.52cm,inner sep=0pt,font=\small},
  square/.style={rectangle,draw,minimum size=0.52cm,inner sep=0pt,font=\small},
  edge/.style={black!80,thick},
  every node/.style={font=\scriptsize}
]
  \node at (1.6,1.0) {trực tiếp};
  \node[round] (a2) at (0,0) {$2$};
  \node[square] (as1) at (1.1,0) {$S_1$};
  \node[square] (as2) at (2.2,0) {$S_2$};
  \node[round] (a8) at (3.3,0) {$8$};
  \draw[edge] (a2)--(as1)--(as2)--(a8);
  \draw[red!70,thick] (as1)--(as2);
  \node[red!70] at (1.65,-0.45) {không hợp lệ};

  \node at (6.1,1.0) {sau khi chèn};
  \node[round] (b2) at (4.4,0) {$2$};
  \node[square] (bs1) at (5.45,0) {$S_1$};
  \node[round] (br) at (6.5,0) {$r$};
  \node[square] (bs2) at (7.55,0) {$S_2$};
  \node[round] (b8) at (8.6,0) {$8$};
  \draw[edge] (b2)--(bs1)--(br)--(bs2)--(b8);
\end{tikzpicture}
\end{center}

\paragraph{UOJ 87.}
Cho cactus, mỗi truy vấn là một tập đỉnh; cần tìm khoảng cách ngắn nhất lớn
nhất giữa hai đỉnh trong tập. Slide ghi ràng buộc tổng quát
\[
  n+\sum |S|\le 300000.
\]
Dựng virtual cây tròn-vuông, đặt trọng số cạnh ảo bằng khoảng cách trên cây
gốc, rồi bài toán trở về dạng truy vấn đường kính trên cấu trúc của BZOJ 2125.

\paragraph{UOJ 189.}
Trong bài \emph{Tài xế tàu ra Tần Xuyên}, cho cactus, hỗ trợ sửa trọng số
cạnh; mỗi truy vấn đưa ra một số đường đi, mỗi đường là đường đơn ngắn nhất
hoặc dài nhất giữa hai đỉnh, và cần tính tổng trọng số của hợp các cạnh được
phủ. Ở đây ``ngắn nhất'' và ``dài nhất'' so theo số cạnh đi qua, không phải
theo trọng số cạnh; đề bảo đảm mọi chu trình đều lẻ. Ràng buộc trong slide là
\[
  n+\sum |S|\le 300000.
\]
Với cây, có thể dùng virtual tree để biến hợp các đường thành tổng các cạnh
được phủ. Với cactus, dựng virtual cây tròn-vuông; trên mỗi chu trình, các
phần bị phủ là hợp của một số đoạn, có thể sắp xếp và cộng bằng tiền tố.

Đối với các cạnh ảo tròn--tròn, chúng tương ứng với hợp của các đường đơn trên
cactus giữa hai đỉnh. Slide tách bốn trường hợp phủ:
\begin{enumerate}
  \item không cạnh nào của chu trình bị phủ;
  \item chỉ phía đường ngắn nhất bị phủ;
  \item chỉ phía đường dài nhất bị phủ;
  \item cả hai phía đều bị phủ.
\end{enumerate}
Nhờ bốn trường hợp này, đóng góp đều có thể biểu diễn bằng một vài tổng độ sâu.

Để hỗ trợ sửa cạnh, có thể duy trì ba loại độ sâu:
\begin{itemize}
  \item \(\operatorname{depTree}\): tổng trọng số các cạnh cây từ gốc;
  \item \(\operatorname{depMin}\): trên mỗi chu trình đi qua, lấy phía ngắn;
  \item \(\operatorname{depMax}\): trên mỗi chu trình đi qua, lấy phía dài.
\end{itemize}
Khi sửa cạnh cây, cộng hiệu vào cả cây con. Khi sửa cạnh trên chu trình lẻ,
tách chu trình tại cha của đỉnh vuông; một phía ảnh hưởng tới
\(\operatorname{depMin}\), phía còn lại ảnh hưởng tới
\(\operatorname{depMax}\). Các cộng đoạn trên cây con có thể dùng cây Fenwick
hoặc cây đoạn.

\subsection{Phân rã và phân tách trên cactus}

Slide dùng bài \emph{UOJ 23: Vua bọ chét xuống Giang Nam} làm ví dụ: cho một
cactus, với mọi \(i\in[1,n)\), hãy in số đường đi đơn độ dài \(i\) xuất phát
từ đỉnh \(1\). Ràng buộc là \(n\le 100000\). Có thể phân rã trọng tâm trên cây
tròn-vuông. Nếu \(g_p(z)\) là hàm sinh của các đường đi từ \(p\) xuống cây con,
thì với một chu trình gồm các đỉnh
\[
  R_1,R_2,\ldots,R_r
\]
không kể gốc chu trình, đóng góp có dạng
\[
  \sum_{i=1}^r g_{R_i}(z)\bigl(z^i+z^{r-i+1}\bigr).
\]
Các phép nhân đa thức ở đỉnh tròn có thể xử lý bằng FFT; ở đỉnh vuông, không
nên nhân thô từng hạng vì bậc có thể lớn. Ta tách thành tổng các đa thức đã
dịch:
\[
  \sum_{q\in S_g}G_q(z)z^{\operatorname{id}_g(q)}
  +\sum_{q\in S_g}G_q(z)z^{r-\operatorname{id}_g(q)+1},
\]
nên chỉ cần cộng đa thức sau khi dịch bậc. Cách này đạt \(O(n\log^2 n)\) trong
phác thảo của slide.

Với bài \emph{UOJ 158: Cactus tĩnh} bản mở rộng dữ liệu, cho một cactus gốc
tại \(1\), mỗi đỉnh đen hoặc trắng, mọi chu trình đều lẻ. Cần hỗ trợ ba thao
tác: đảo màu mọi đỉnh trên đường ngắn nhất từ một đỉnh tới gốc; đảo màu mọi
đỉnh trên đường đơn dài nhất từ một đỉnh tới gốc; và hỏi trong sub-cactus của
một đỉnh có bao nhiêu đỉnh đen. Ràng buộc là \(n,q\le 200000\).

Ta có thể mô phỏng ý tưởng heavy-light decomposition trên cây tròn-vuông.
Sub-cactus tương ứng với cây con, nên truy vấn đếm trong cây con giống bài
cây. Phần khó là cập nhật đường ngắn nhất hoặc dài nhất trên một chuỗi nặng.
Ta phân loại các đỉnh tròn trên chuỗi: đỉnh bắt buộc đi qua, đỉnh chỉ nằm trên
đường ngắn nhất, và đỉnh chỉ nằm trên đường dài nhất. Mỗi loại được duy trì
riêng bằng cây đoạn trên một thứ tự DFS được chỉnh lại cho đỉnh vuông: trước
tiên đưa tất cả đỉnh tròn kề đỉnh vuông vào thứ tự, rồi mới đệ quy vào các cây
con của chúng. Hình ví dụ trong slide có thứ tự DFS
\[
  \text{gốc},2,0,1,5,4,3,6,7,9,8,
\]
minh họa rằng các đỉnh tròn kề một đỉnh vuông được đặt liên tiếp trước khi đi
sâu vào các cây con.

\section{So sánh với cây co chu trình}

Cây co chu trình và cây tròn-vuông đều biến cactus thành cây, và đều có thể
dùng cho phân tách hoặc virtual tree. Khác biệt chính là:
\begin{itemize}
  \item cây co chu trình là cây có gốc; cây tròn-vuông tự nhiên là cây không
        gốc;
  \item cây tròn-vuông giữ tương ứng một-một giữa đỉnh gốc và đỉnh tròn, trong
        khi cây co chu trình có thể làm mất thông tin vị trí trong chu trình;
  \item mọi cây tròn-vuông hợp lệ đều tương ứng với một cactus, còn nhiều
        cactus khác nhau có thể cho cùng một cây co chu trình;
  \item khi cần biết một đỉnh thuộc sub-cactus nào của một chu trình, cây
        tròn-vuông có thể dùng nhảy trên cây hoặc HLD trực tiếp.
\end{itemize}
Vì không co các đỉnh gốc, cây tròn-vuông thường có tính chất tốt hơn cho
virtual cactus, phân rã và các bài cần bảo toàn vị trí trên chu trình.

\section{Cây tròn-vuông tổng quát}

Với một đồ thị vô hướng tổng quát, nếu ta tạo một đỉnh vuông cho mỗi thành
phần song liên thông đỉnh rồi nối nó với mọi đỉnh trong thành phần đó, ta nhận
được một cấu trúc tương tự, gọi là \term{cây tròn-vuông tổng quát}. Theo định
lý tương đương ở trên, có thể xem mỗi thành phần song liên thông như một chu
trình của một cactus tương đương, dù thứ tự trên ``chu trình'' là tùy chọn.
Cách nhìn này giúp biến một số bài trên thành phần song liên thông thành bài
trên cây.

\paragraph{Bài ``Thương nhân''.}
Cho đồ thị có trọng số đỉnh. Mỗi truy vấn hỏi trên tất cả đường đi đơn giữa
hai đỉnh, giá trị nhỏ nhất có thể gặp là bao nhiêu. Dựng cây tròn-vuông tổng
quát và đặt trọng số của mỗi đỉnh vuông bằng trọng số nhỏ nhất của các đỉnh
tròn kề nó. Khi đó truy vấn trở thành lấy min trên đường đi trong cây.

\paragraph{CodeChef SADPAIRS.}
Cho đồ thị có màu đỉnh; với mỗi đỉnh, hỏi nếu xóa đỉnh đó thì có bao nhiêu cặp
đỉnh cùng màu bị tách rời. Xét riêng từng màu, dựng virtual tree trên cây
tròn-vuông tổng quát. Mỗi cạnh ảo tương ứng với một đường trên cây thật và
đóng góp
\[
  \text{số đỉnh trong cây con}\times \text{số đỉnh ngoài cây con};
\]
có thể dùng hiệu trên cây để cộng đóng góp. Với LCA tuyến tính bằng RMQ và sắp
xếp virtual tree bằng radix sort, tổng độ phức tạp có thể đạt tuyến tính theo
kích thước đầu vào.

\section*{Kết luận}

Các kỹ thuật trong tài liệu đều cùng một mục tiêu: đưa phần ``không phải cây''
của đồ thị về một số nút hoặc trạng thái nhỏ, rồi tận dụng các thuật toán quen
thuộc trên cây. Cây DFS thích hợp cho đồ thị gần cây và các ràng buộc do cạnh
ngoài cây phủ đường đi. Cây BFS thích hợp cho phân lớp. Cây tròn-vuông và biến
thể tổng quát là công cụ mạnh cho cactus và thành phần song liên thông, đặc
biệt khi cần virtual tree, truy vấn đường đi, phân rã hoặc heavy-light
decomposition.

Các slide tổng kết nhấn mạnh ba tác dụng của cây tròn-vuông: nó là cách tốt
hơn và tổng quát hơn để xử lý nhiều bài cactus, làm giảm độ khó khi biến cactus
về cây; nó liên hệ cactus với thành phần song liên thông đỉnh, khiến cách nghĩ
cụ thể hơn; và nó mở đường cho nhiều bài cactus hoặc điểm song liên thông hơn.
Phần triển vọng mong người đọc nghĩ ra thêm các ý tưởng bài cactus tốt hơn và
phát triển cây tròn-vuông tới các bài cactus động.

Phần cảm ơn của nguồn ghi lời cảm ơn cha mẹ, thầy cô và các huấn luyện viên
Chen Ying, Yu Linyun, Chen Xumin, Zhang Ruizhe, cùng đàn anh Dong Kefan; đồng
thời cảm ơn Liu Yifan (LGG) của Fuzhou No. 1 High School đã cùng tác giả nghĩ
ra tên ``cây tròn-vuông'', bổ sung chi tiết, và kiểm tra mã khi tác giả viết
\emph{Cactus tĩnh} và \emph{Tài xế tàu ra Tần Xuyên}.

\end{document}
